\documentclass[10pt,a4paper]{amsart}
\usepackage[margin=19mm]{geometry}
\usepackage{amsmath,amssymb,mathtools}
\usepackage[scr=rsfs,cal=euler]{mathalfa}
\usepackage{xfrac}
\usepackage[unicode,hidelinks,bookmarks=false]{hyperref}
\newcommand{\ie}{i.e.\ }
\newcommand{\cf}{cf.\ }
\newcommand{\resp}{respectively\ }
\newcommand{\Subsection}{\subsection}
\providecommand{\nobreakdash}{\nobreak\hbox{-}}
\setlength{\emergencystretch}{2em}
\multlinegap0pt
\usepackage{bm}
\usepackage{bbm}
\usepackage{dsfont}
\usepackage{xspace}
\usepackage{cancel}
\usepackage{mathtools}
\usepackage{amsthm}
\makeatletter
\@ifundefined{tw}{\newtheorem{tw}{Tw}}{}
\@ifundefined{prop}{\newtheorem{prop}[tw]{Proposition}}{}
\makeatother
\title[A Study of the Spectral Sequence for Locally Free Isometric ]{A Study of the Spectral Sequence for Locally Free Isometric Actions of Abelian Lie Groups}
\author{Pawe{\l} Ra\'zny}
\date{Source: arXiv:2506.12422v2 (CC BY 4.0)}
\begin{document}
\maketitle

\noindent\emph{Source and licence.} A Study of the Spectral Sequence for Locally Free Isometric Actions of Abelian Lie Groups. Version arXiv:2506.12422v2 (CC BY 4.0), distributed under the Creative Commons Attribution 4.0 International licence, as recorded in the arXiv licence metadata for this exact version and shown on the abstract page https://arxiv.org/abs/2506.12422v2. Full rights notice: licenses/arxiv-2506.12422-CC-BY-4.0.txt.\par\smallskip

\noindent\textit{Editorial scope.} Counted unit: the Prop whose source label is \texttt{InvFor} in the article (source0.tex lines 176--181, source label \texttt{InvFor}), delivered together with the article's own complete original proof of it (source0.tex lines 182--202, one \texttt{proof} environment of 21 lines). No mathematics is written, repaired or re-derived: statement and proof are the article's own text line for line. Required context retained uncounted: none. Every definition, display and label of those lines is retained unchanged. Omitted from this package and not redistributed: 1--175, 203--592. Nothing inside a retained excerpt is omitted or altered: every retained body is delivered exactly as the article's lines are written, so no omission receipt of any kind accompanies this package. Citations: the retained excerpts carry no citation command at all, so no bibliography entry of the article is transcribed and no reference is redistributed. Reference adaptation: none was needed and none was made; every reference inside the delivered unit resolves inside it, so no unresolved reference or citation is produced. Printed slips kept literal: no printed word, symbol, hypothesis or number of the counted statement or of its proof was altered. Layout and class: the article's own class is \texttt{sigma} with packages including ifpdf, tikz-cd, mathtools; the wrapper is the lane's pinned \texttt{amsart} editorial wrapper at 10pt on A4, which restates the theorem-like environments the retained text needs and adds the article's own macro definitions that the retained text needs (none), with extra wrapper packages (bm, bbm, dsfont, xspace, cancel, mathtools). The delivered numbering is the unit's own. Assets: no retained span contains a figure environment, a graphics-inclusion command, a PDF-page inclusion, a table or a TikZ picture; no image file is copied into this package, the package declares no asset dependency and no image file is redistributed with it. Licence: version arXiv:2506.12422v2 of this article is distributed under the Creative Commons Attribution 4.0 International licence, as recorded in the arXiv licence metadata for this exact version and shown on the abstract page https://arxiv.org/abs/2506.12422v2. A full rights notice is delivered at licenses/arxiv-2506.12422-CC-BY-4.0.txt. Characters: every retained excerpt is pure ASCII, so the delivered engine, XeLaTeX with its default Latin Modern text font, reports no missing character.
\begin{prop}\label{InvFor} Let $(M^{n+s},g)$ be a compact Riemannian manifold with a locally free action of a connected abelian Lie group $G\subset \operatorname{Diff}(M)$ by isometries. Then the forms ${\rm d}\eta_i$ are basic. Moreover, the following conditions are equivalent:
\begin{enumerate}\itemsep=0pt
\item[$(1)$] $\alpha$ is a $\overline{G}$-invariant form on $M$.
\item[$(2)$] \smash{$\alpha=\alpha_0+\sum_{k=1}^p\sum_{1\leq i_1<\cdots <i_k\leq s}\eta_{i_1}\cdots \eta_{i_k}\alpha_{i_1,\dots,i_k}$}, where $\alpha_0$ and $\alpha_{i_1,\dots,i_k}$ are basic for all indices $1\leq i_1<\cdots <i_k\leq s$.
\end{enumerate}
\end{prop}

\begin{proof}
Firstly, let us show that the forms ${\rm d}\eta_i$ are indeed basic. To see this, let us note that since $\overline{G}$ acts by isometries and the flows of $\xi_i$ preserve $T\mathcal{F}$, they have to preserve $T\mathcal{F}^{\perp}$ as well. Consider the formula
\[{\rm d}\eta_i(X_0,X_1)= X_0.(\eta_i(X_1))-X_1.(\eta_i(X_0)) -\eta_i([X_0,X_1]).\]
To prove that this form indeed vanishes on $T\mathcal{F}$, it suffices to consider $X_0$ to be an $\mathbb{R}$-linear combination of $\xi_1,\dots,\xi_p$ (by tensoriality). In this case, $\eta_i(X_0)$ is constant and hence the second term vanishes. On the other hand, we can without loss of generality (by tensoriality) take $X_1$ to be a sum of an $\mathbb{R}$-linear combination of $\xi_1,\dots,\xi_p$ and a vector field in $T\mathcal{F}^{\perp}$. Hence, similarly as before, the first term in the sum vanishes. Moreover, the final term also vanishes since $[X_0,X_1]$ has to be a section of $T\mathcal{F}^{\perp}$ (since the brackets of the form $[\xi_i,\xi_j]$ vanish and $T\mathcal{F}^{\perp}$ is preserved by the flows of $\xi_i$ as already remarked). Finally, for any $X\in\Gamma(T\mathcal{F})$, we have{\samepage
\[\mathcal{L}_X{\rm d}\eta_j={\rm d}\iota_X{\rm d}\eta_j+\iota_X{\rm d}{\rm d}\eta_j=0,\]
which implies that ${\rm d}\eta_j$ is indeed basic.}

Now assume that the second condition is true. Then it can be easily computed that for any~$\xi_j$ the equality $\mathcal{L}_{\xi_j}\alpha=0$ holds. Which in turn implies that $\alpha$ is $G$-invariant and consequently $\overline{G}$-invariant.

Now let us write the invariant form $\alpha$ as
\[\alpha=\alpha_0+\sum\limits_{k=1}^s\sum\limits_{1\leq i_1<\cdots <i_k\leq s}\eta_{i_1}\cdots \eta_{i_k}\alpha_{i_1,\dots,i_k},\]
where $\alpha_{i_1,\dots,i_k}$ are transverse for all indices $1\leq i_1<\cdots <i_k\leq s$. Due to the well-known formula
\[\mathcal{L}_{X}i_{Y}-i_{Y}\mathcal{L}_{X}=i_{[X,Y]},\]
we get that $i_{\xi_i}$ and $\mathcal{L}_{\xi_j}$ commute for $i,j\in\{1,\dots,s\}$. We shall now prove that the forms~$\alpha_0$ and~${\alpha_{i_1,\dots,i_k}}$ are basic by reverse induction on the number of indices. Hence, we start by proving that $\alpha_{1,\dots,s}$ is basic. Since $\alpha$ is basic and the vector fields $\xi_i$ are Killing, we have for any ${i\in\{1,\dots,s\}}$ the following equalities:
\[0=\mathcal{L}_{\xi_i}\alpha=i_{\xi_s}i_{\xi_{s-1}}\cdots i_{\xi_1}\mathcal{L}_{\xi_i}\alpha=\mathcal{L}_{\xi_i}i_{\xi_{s}}i_{\xi_{s-1}}\cdots i_{\xi_1}\alpha=\mathcal{L}_{\xi_i}\alpha_{1,\dots,s},\]
which proves that $\alpha_{1,\dots,s}$ is basic.

For the induction step, let us assume that all the $\alpha_{i_1,\dots,i_{k}}$ for $p\geq k>K$ are basic. We shall show that all $\alpha_{i_1,\dots,i_{K}}$ are basic as well. Using the assumption, we get for any $i\in\{1,\dots,s\}$ the following equalities:
\[0=\mathcal{L}_{\xi_i}\alpha=i_{\xi_{i_K}}i_{\xi_{i_{K-1}}}\cdots i_{\xi_{i_1}}\mathcal{L}_{\xi_i}\alpha=\mathcal{L}_{\xi_i}i_{\xi_{i_K}}i_{\xi_{i_{K-1}}}\cdots i_{\xi_{i_1}}\alpha=\mathcal{L}_{\xi_i}\alpha_{i_1,\dots,i_K},\]
which proves that $\alpha_{i_1,\dots,i_K}$ are basic for any set of indices $1\leq i_1<\cdots <i_k\leq s$.
\end{proof}

\end{document}
